%titolo{Scrolla le ali}
%autore{}
%album{}
%tonalita{DO}
%famiglia{Scout}
%gruppo{}
%momenti{}
%identificatore{scrolla_le_ali}
%data_revisione{2026_07_01}
%trascrittore{Modanese Tobia}

\beginsong{Scrolla le ali}

\beginchorus
Scrolla le \[C]ali, cocci\[F]nella, 
e con\[C]trolla un po' l'an\[G]tenna, 
scrolla le \[C]ali che si \[G]parte, 
un gran \[F]volo si \[G]farà! 
Prato, \[F]bosco e poi mon\[C]tagna, 
trove\[F]remo \[G]l'aqui\[C]la. 
\endchorus

\beginverse
\[C]Sia\[C7]mo \[F]otto cocci\[G]nelle 
e vo\[C]lia\[C7]mo \[F]libere e so\[G]relle. 
\[A-]Guar\[D-]da, \[F]siamo sopra un \[G]prato, 
gli \[B&]animali ci \[F]hanno salu\[G]tato. 
\endverse
\chordsoff
\beginverse
La cicala sta nel prato
e ci da un consiglio spensierato:
chi va forte lascia soli gli altri
chi va piano sempre solo volerà.
\endverse

\beginverse
Siamo nel bosco più vicino
e parliamo col vecchio porcospino
poi l'inverno comincia ad arrivare
un bel pino ci lascia riparare.
\endverse

\beginverse
Che fatica siamo già in montagna
già si sente qualcuno che si lagna
una capra ci invita a continuare
sulla vetta potremo festeggiare.
\endverse

\beginverse
Ma che bello il viaggio è completato
qui dal monte vediamo il bosco e il prato
con Arcanda possiamo chiacchierare
le avventure vogliamo continuare.
\endverse

\endsong
